%!TEX root = ../thesis.tex
%*******************************************************************************
%************************** Materials and Methods ******************************
%*******************************************************************************

\chapter{Materials and Methods}
\label{ch:methods}

\ifpdf
    \graphicspath{{Chapter2/Figs/Raster/}{Chapter2/Figs/PDF/}{Chapter2/Figs/}}
\else
    \graphicspath{{Chapter2/Figs/Vector/}{Chapter2/Figs/}}
\fi

%********************************** Section 2.1 ********************************
\section{System Architecture Overview}

The recirculation platform comprises four principal subsystems integrated around an ESP32-WROOM-32 microcontroller: (i) a Bartels MP6 piezoelectric micropump with high-voltage driver and DAC, (ii) a Sensirion SLF3S liquid flow sensor, (iii) a Bartels mp-bt bubble trap, and (iv) a host PC running a Python-based GUI. Figure~\ref{fig:system_architecture} illustrates the complete system architecture.

% TODO: Add system architecture figure when ready
%\begin{figure}[htbp]
%    \centering
%    \includegraphics[width=0.9\textwidth]{system_architecture}
%    \caption{System architecture of the biomimetic recirculation platform. The ESP32 microcontroller coordinates pump drive electronics, sensor acquisition, and closed-loop PID control, communicating with the host PC via UART for real-time monitoring and data recording.}
%    \label{fig:system_architecture}
%\end{figure}

The fluid circuit forms a closed recirculation loop: liquid is drawn from a temperature-controlled reservoir by the micropump, passes through the bubble trap, flows through the SLF3S sensor, enters the microfluidic test section, and returns to the reservoir. All fluidic connections use standard \SI{1.6}{\milli\metre} inner diameter tubing. This system advances the first-generation ``Artery in a Box'' prototype \citep{Kacerik2024} by introducing closed-loop PID regulation, real-time shear rate computation, and comprehensive data logging.

%********************************** Section 2.2 ********************************
\section{Micropump and Drive Electronics}

The Bartels MP6 is a piezoelectric diaphragm micropump that generates flow through periodic deformation of a lead zirconate titanate (PZT) membrane bonded to a polyphenylsulfone (PPSU) body \citep{BartelsMP6}. Two passive check valves ensure unidirectional flow. Key specifications include a maximum flow rate of \SI{9}{\milli\litre\per\minute}, maximum back pressure of \SI{450}{\milli\bar}, internal dead volume of approximately \SI{30}{\micro\litre}, and power consumption of approximately \SI{50}{\milli\watt}.

Flow rate is governed by the drive signal \textbf{amplitude} (\SIrange{85}{250}{\volt} peak-to-peak) and \textbf{frequency} (\SIrange{25}{226}{\hertz}). In this system, the frequency is held constant at the optimal operating point while the amplitude is varied dynamically by the PID controller, as amplitude provides a more linear control actuation path.

The MP6 requires a dedicated high-voltage driver board (Bartels mp-Driver) accepting three inputs from the ESP32: an enable signal, a clock signal defining drive frequency via PWM, and an analogue voltage setting the amplitude. The amplitude is set by a Microchip MCP4726 12-bit DAC (I\textsuperscript{2}C, address \texttt{0x61}), whose output voltage is:

\begin{equation}
    V_{\text{out}} = \frac{\text{DAC code}}{4096} \times V_{\text{DD}}
    \label{eq:dac}
\end{equation}

During calibration, the actual supply voltage was measured as \SI{4.734}{\volt} rather than the nominal \SI{5.0}{\volt}, requiring correction of the amplitude-to-voltage mapping. The corrected range spans \SI{0.35}{\volt} (code 287) to \SI{1.30}{\volt} (code 1065).

%********************************** Section 2.3 ********************************
\section{Flow Sensor: Sensirion SLF3S}

Real-time flow measurement is provided by a Sensirion SLF3S liquid flow sensor operating on the calorimetric (thermal mass flow) principle \citep{SensirionSLF3S}. A micro-heater establishes a thermal gradient; two symmetrically placed thermopile sensors detect its displacement by the flow, producing a signal proportional to mass flow rate.

The SLF3S-0600F variant was selected for its measurement range of \SIrange{0}{2000}{\micro\litre\per\minute}, optimised for microfluidic applications. The sensor communicates via I\textsuperscript{2}C at address \texttt{0x08}, providing 16-bit digital flow readings with an accuracy of $\pm$5\% of measured value or \SI{0.5}{\micro\litre\per\minute}. A sampling rate of \SI{10}{\hertz} was selected for the control loop. The built-in \textbf{air-in-line detection} capability flags the presence of bubbles, enabling real-time alerts critical for protecting biological samples.

%********************************** Section 2.4 ********************************
\section{Key Fluidic Parameters}

For a Newtonian fluid in fully developed laminar flow through a circular tube, the wall shear rate $\dot{\gamma}_w$ is:

\begin{equation}
    \dot{\gamma}_w = \frac{4Q}{\pi r^3}
    \label{eq:shear_rate}
\end{equation}

\noindent where $Q$ is the volumetric flow rate and $r$ is the tube inner radius. The corresponding mean flow velocity is $\bar{v} = Q / (\pi r^2)$. The assumption of laminar flow is verified by the Reynolds number: for water at \SI{37}{\celsius} flowing at \SI{1}{\milli\litre\per\minute} through \SI{1.6}{\milli\metre} tubing, $Re \approx 14$, well within the laminar regime.

Table~\ref{tab:blood_vessels} summarises the typical shear rates encountered in human blood vessels \citep{Popel2005, Sakariassen2015}. The system's flow rate range enables coverage of shear rates spanning venous through arteriolar conditions.

\begin{table}[htbp]
    \centering
    \caption{Typical wall shear rates in human blood vessels \citep{Popel2005, Sakariassen2015}.}
    \label{tab:blood_vessels}
    \begin{tabular}{@{}lcc@{}}
        \toprule
        Vessel type & Diameter (\si{\milli\metre}) & Wall shear rate (\si{\per\second}) \\
        \midrule
        Large veins       & 5--10   & 20--200    \\
        Small veins       & 1--5    & 200--500   \\
        Arteries          & 3--6    & 300--800   \\
        Arterioles        & 0.01--0.3 & 500--5000  \\
        \bottomrule
    \end{tabular}
\end{table}

%********************************** Section 2.5 ********************************
\section{PID Closed-Loop Control}

The closed-loop system employs a discrete PID algorithm to regulate flow rate by adjusting pump drive amplitude:

\begin{equation}
    u(k) = K_p \, e(k) + K_i \sum_{j=0}^{k} e(j) \, \Delta t + K_d \, \frac{e(k) - e(k-1)}{\Delta t}
    \label{eq:pid}
\end{equation}

\noindent where $e(k) = Q_{\text{target}} - Q_{\text{measured}}$ is the flow rate error and $\Delta t = \SI{100}{\milli\second}$ is the sampling interval. The controller output is clamped to the valid amplitude range (80--250), and the integral accumulator is bounded to $\pm 500$ to prevent windup. Default gains ($K_p = 1.0$, $K_i = 0.1$, $K_d = 0.01$) were determined through manual tuning, prioritising stable tracking with minimal overshoot. A safety mechanism triggers a \texttt{FLOW\_ERR} alert if measured flow deviates from the target by more than 20\% for \SI{10}{\second}.

%********************************** Section 2.6 ********************************
\section{Fluidic Circuit and Bubble Management}

The complete circuit forms a closed recirculation loop with liquid maintained at \SI{37}{\celsius} using a hot plate. The Bartels mp-bt bubble trap, positioned between pump and sensor, employs a hydrophobic microporous membrane (internal volume $\sim$\SI{100}{\micro\litre}) that passively separates gas from the liquid stream, requiring \SIrange{100}{200}{\milli\bar} of back pressure. Before each experiment, priming at maximum pump settings flushes the circuit until the flow sensor reading stabilises and no further air-in-line alerts are triggered.
